\documentclass[11pt]{article}

\usepackage[margin=1in]{geometry}
\usepackage{amsmath,amssymb,amsthm,microtype,needspace}
\usepackage[T1]{fontenc}
\usepackage[hidelinks]{hyperref}

\newtheorem{theorem}{Theorem}[section]
\newtheorem{proposition}[theorem]{Proposition}
\newtheorem{lemma}[theorem]{Lemma}
\newcommand{\A}{\mathbb A}
\newcommand{\Pone}{\mathbb P^1}
\makeatletter
\renewcommand{\@seccntformat}[1]{\csname the#1\endcsname\hspace{0.75em}}
\makeatother

\title{Generic Discriminants of Polynomial Tangent Pencils}
\author{Roy van Rijn\\{\small Independent researcher}}
\date{17 November 2026}
\hypersetup{
  pdftitle={Generic Discriminants of Polynomial Tangent Pencils},
  pdfauthor={Roy van Rijn},
  pdfsubject={Singularities of the discriminant of H(W)-sW+t},
  pdfkeywords={polynomial pencil, discriminant curve, cusp, node,
    tangent line}
}

\begin{document}
\maketitle

\begin{abstract}
Let \(H\) be a polynomial of degree \(n\geq3\) over an algebraically
closed field of characteristic zero.  The repeated-root locus of the pencil
\[
  H(W)-sW+t
\]
is parameterized by
\[
  r\longmapsto\bigl(H'(r),rH'(r)-H(r)\bigr).
\]
We give a self-contained contact-stratum proof of this classical dual-curve
specialization: for \(H\) in a nonempty Zariski-open set, its projective
closure is a rational plane curve of degree \(n\), smooth at its unique
point at infinity, with exactly \(n-2\) ordinary cusps and
\(\frac12(n-2)(n-3)\) ordinary nodes, and no other singularities.
The proof uses the contact of the graph of \(H\) with affine lines to
exclude every nongeneric singularity before the genus formula is applied.
\end{abstract}

\section{The theorem}

Let \(k\) be an algebraically closed field of characteristic zero and write
\[
  H(W)=\sum_{j=2}^{n}h_jW^j,\qquad h_n\neq0.              \tag{1.1}
\]
Discarding the affine-linear part of \(H\) only changes coordinates in the
\((s,t)\)-plane, so (1.1) entails no loss of generality.  Put
\[
  E_{s,t}(W)=H(W)-sW+t.
\]
A repeated root \(r\) satisfies
\[
  s=H'(r),\qquad t=rH'(r)-H(r).                           \tag{1.2}
\]
Thus the reduced discriminant curve is the image of
\[
  \nu_H:\A^1\longrightarrow\A^2,\qquad
  r\longmapsto\bigl(H'(r),rH'(r)-H(r)\bigr).              \tag{1.3}
\]
Geometrically, \((s,-t)\) records the affine line of slope \(s\) tangent to
the graph \(Y=H(W)\) at \(r\).  Cusps arise from flexes of the graph and
nodes from lines tangent at two distinct points.  This is a special
one-parameter instance of the classical geometry of dual plane
curves.  The contribution of this note is an elementary verification of
genericity and the behavior at infinity in the thin family of polynomial
graphs, followed by the usual genus count.  No new general Pl\"ucker formula
or priority for the numerical counts is claimed.

For comparison, Yuran's Newton-polygon theorem gives only ordinary contacts
under the sufficient hypothesis that the polygon contains \(5\Delta\)
\cite[Theorems~1.3 and~1.16]{Yuran}.  The graph \(Y-H(W)=0\) has Newton
polygon contained in \(0\leq\deg_Y\leq1\), so that hypothesis does not
apply to this family.  Yuran's conditional counting theorems assume the
corresponding contact exclusions; we verify the needed exclusions directly
by the three factorization strata below.  This comparison identifies the
proof task here without asserting that these particular counts were
previously unknown.

\Needspace{16\baselineskip}
\begin{theorem}\label{thm:main}
For every \(n\geq3\), there is a nonempty Zariski-open subset of the
degree-\(n\) polynomials such that the projective discriminant curve of
\(H(W)-sW+t\) has:
\begin{enumerate}
\item degree \(n\) and normalization \(\Pone\);
\item one point at infinity, which is smooth;
\item exactly \(n-2\) ordinary cusps;
\item exactly \(\frac12(n-2)(n-3)\) ordinary nodes;
\item no other singularities.
\end{enumerate}
\end{theorem}

The result concerns the reduced discriminant.  No assertion is made here
about the multiplicity with which a resultant defines it.

For the cubic \(H(W)=W^3\), the parametrization is
\(\nu_H(r)=(3r^2,2r^3)\).  Its affine equation is
\(27t^2=4s^3\), so the origin is its unique ordinary cusp and there are no
nodes.  Homogenizing gives \(27T^2Z=4S^3\), whose sole point at infinity
\([0:1:0]\) is smooth since the derivative with respect to \(Z\) is nonzero.
This is the \(n=3\) case of all five conclusions.

\section{Normalization and infinity}

Homogenizing (1.3) gives
\[
  \bar\nu_H:\Pone\longrightarrow\mathbb P^2,\qquad
  [R:U]\longmapsto[S_H:T_H:Z],                            \tag{2.1}
\]
where
\begin{align*}
 S_H&=\sum_{j=2}^n jh_jR^{j-1}U^{n-j+1},\\
 T_H&=\sum_{j=2}^n (j-1)h_jR^jU^{n-j},\\
 Z&=U^n.
\end{align*}
These forms have degree \(n\) and no common zero.

\begin{proposition}\label{prop:normalization}
For every polynomial \(H\) of exact degree \(n\), the map
\(\bar\nu_H\) is finite and birational onto an integral plane curve of
degree \(n\).  Its image has a unique point at infinity, and that point is
smooth.
\end{proposition}

\begin{proof}
On the affine chart,
\[
  \frac{ds}{dr}=H''(r),\qquad
  \frac{dt}{dr}=rH''(r).
\]
At the generic point of the image this recovers \(r=dt/ds\), so
\(\bar\nu_H\) is birational.  A nonconstant morphism from a proper curve to
its image is finite.  Since (2.1) is given by base-point-free forms of
degree \(n\), birationality makes the image a curve of degree \(n\).

The equation \(Z=0\) has the unique preimage \([1:0]\), whose image is
\([0:1:0]\).  Put \(q=U/R\) and work in the target chart \(T_H\neq0\).
Then
\[
  \frac{S_H}{T_H}=\frac{n}{n-1}q+O(q^2),\qquad
  \frac{Z}{T_H}=\frac{q^n}{(n-1)h_n}+O(q^{n+1}).          \tag{2.2}
\]
The first series has nonzero linear coefficient.  It is therefore a local
parameter on the image, which is formally the graph of the second series as
a function of the first.  The point at infinity is smooth.
\end{proof}

\section{Contact strata}

Let \(\ell_r\) be the tangent line to the graph of \(H\) at \(r\), and put
\[
  c(r)=\operatorname{ord}_{W=r}\bigl(H(W)-\ell_r(W)\bigr).
\]
One always has \(c(r)\geq2\).  Differentiating (1.3) gives
\[
  \nu_H'(r)=H''(r)(1,r).                                 \tag{3.1}
\]
Thus \(c(r)=2\) gives a smooth immersed branch.  If \(c(r)=3\), then
\(H'''(r)\neq0\), and
\[
  \det\bigl(\nu_H''(r),\nu_H'''(r)\bigr)
  =2H'''(r)^2\neq0.                                      \tag{3.2}
\]
The image branch is therefore an ordinary cusp.  Contact of order at least
four is the first bad one-branch event.

If two distinct parameters \(r\) and \(u\) have the same image and both
branches are unramified, their tangent directions are \((1,r)\) and
\((1,u)\).  Their determinant is \(u-r\neq0\), so the two branches meet in
an ordinary node.

It remains to show that every other event is nongeneric.  For a contact
pattern \(\mu=(\mu_1,\ldots,\mu_k)\), write
\[
  M_\mu(W)=\prod_{i=1}^k(W-r_i)^{\mu_i},\qquad
  m=\sum_i\mu_i.
\]
An affine line having the prescribed contacts with the graph of \(H\)
gives a factorization
\[
  H(W)-\ell(W)=M_\mu(W)Q(W),\qquad \deg Q\leq n-m.        \tag{3.3}
\]
After affine-linear terms are discarded, such polynomials vary in a family
of dimension at most
\[
  k+(n-m+1).                                             \tag{3.4}
\]
The coefficient space in (1.1) has dimension \(n-1\).

\begin{lemma}\label{lem:bad}
Each of the contact patterns
\[
  (4),\qquad(3,2),\qquad(2,2,2)                          \tag{3.5}
\]
is contained in a proper closed subset of the degree-\(n\) coefficient
space.
\end{lemma}

\begin{proof}
For each pattern in (3.5), the dimension in (3.4) is \(n-2\), whenever the
pattern is possible.  The closure of the image of the factorization family
(3.3) therefore has dimension at most \(n-2\), hence is proper in the
\((n-1)\)-dimensional coefficient space.  Taking the closure also includes
collisions among the marked contact points and degree drops in \(Q\).
\end{proof}

\begin{lemma}\label{lem:exhaustion}
Outside the three closed loci in Lemma~\ref{lem:bad}, every affine
singularity of the discriminant is an ordinary cusp or an ordinary node.
\end{lemma}

\begin{proof}
A nonordinary single branch has contact at least four and belongs to the
\((4)\)-locus.  If an image point has two normalization preimages and one
branch is ramified, its contact pattern dominates \((3,2)\).  If it has at
least three preimages, any three give the pattern \((2,2,2)\).  The only
remaining singular image has exactly two unramified preimages, and (3.1)
makes their tangent directions distinct.
\end{proof}

Proposition~\ref{prop:normalization} shows uniformly that no additional
singularity occurs at infinity.  The complement of the three bad loci is
therefore a nonempty Zariski-open set on which the discriminant has only
ordinary cusps and nodes.

\section{The count}

On the open set just constructed, \(H''\) is squarefree.  Its \(n-2\) roots
are precisely the ramification points of (1.3), and (3.2) shows that their
images are \(n-2\) distinct ordinary cusps.

The discriminant is an integral plane curve of degree \(n\) with
normalization \(\Pone\).  Its arithmetic genus is
\[
  p_a=\frac{(n-1)(n-2)}2,
\]
whereas its geometric genus is zero.  Ordinary nodes and ordinary cusps
both have delta invariant one.  If \(N\) is the number of nodes, the genus
formula gives
\[
  \frac{(n-1)(n-2)}2=(n-2)+N.
\]
Hence
\[
  N=\frac{(n-2)(n-3)}2,
\]
which proves Theorem~\ref{thm:main}.

\section*{Formalization scope}

The Lean development checks, without additional axioms, the
repeated-root parametrization, the contact and jet calculations, the formal
local analysis at infinity, the complete affine normalization-fiber
exhaustion under the three explicit bad-contact exclusions, and the distinct
cusp and numerical counts.
The global dimension-of-image argument, the normalization and degree
theorem for the parametrized image, the scheme-theoretic cusp and node
criteria, and the genus--delta formula are supplied by the written proof
and lie outside this partial formalization.  A versioned source snapshot
with its pinned toolchain is available as~\cite{FormalSources}; build and
scope instructions accompany the manuscript supplement.  This availability
statement does not report a new Lean replay.

\paragraph{AI assistance and review.}
ChatGPT (OpenAI) assisted with mathematical arguments, literature checks,
exposition, proof review and verification code.

\begin{thebibliography}{9}

\bibitem{FormalSources}
R.~van Rijn, \emph{Polynomial tangent pencils in Lean},
repository edition \texttt{c32b9b4}, 17 November 2026,
\href{https://github.com/royvanrijn/jacobian-research/tree/c32b9b4d273cf7988e85184a0d5eae9c488727fd/research/foundations/formal/discriminant-pencils}{versioned partial formalization}.

\bibitem{Yuran}
A.~Yuran,
\emph{Pl\"ucker formulas for plane algebraic curves with a given Newton
polygon},
arXiv:2204.04626.

\end{thebibliography}

\end{document}
